\textbf{Phần 1. Trắc nghiệm lựa chọn}
\setcounterpageref{ex}{0}

\begin{ex} %[5D6N34]
Đổi $9{,}02\text{ km}$ ra mét ta được kết quả là
\choice
{$902\text{ m}$}
{$90\,200\text{ m}$}
{\True $9\,020\text{ m}$}
{$9{,}020\text{ m}$}
\loigiai{
Vì $1\text{ km} = 1\,000\text{ m}$ nên $9{,}02\text{ km} = 9{,}02 \times 1\,000 = 9\,020\text{ m}$.\\
\textbf{Đáp án của câu hỏi là C}
}
\end{ex}

\begin{ex} %[5D6N34]
Đổi $6{,}35\text{ tấn}$ ra ki-lô-gam ta được kết quả là
\choice
{$635\text{ kg}$}
{\True $6\,350\text{ kg}$}
{$63\,500\text{ kg}$}
{$6{,}350\text{ kg}$}
\loigiai{
Vì $1\text{ tấn} = 1\,000\text{ kg}$ nên $6{,}35\text{ tấn} = 6{,}35 \times 1\,000 = 6\,350\text{ kg}$.\\
\textbf{Đáp án của câu hỏi là B}
}
\end{ex}

\begin{ex} %[5D6H34]
Diện tích của Hồ Gươm khoảng $0{,}12\text{ km}^2$. Diện tích Hồ Gươm tính theo héc-ta và mét vuông lần lượt là
\choice
{\True $12\text{ ha}$ và $120\,000\text{ m}^2$}
{$1{,}2\text{ ha}$ và $12\,000\text{ m}^2$}
{$120\text{ ha}$ và $1\,200\,000\text{ m}^2$}
{$12\text{ ha}$ và $12\,000\text{ m}^2$}
\loigiai{
Vì $1\text{ km}^2 = 100\text{ ha} = 1\,000\,000\text{ m}^2$ nên:\\
$0{,}12\text{ km}^2 = 0{,}12 \times 100 = 12\text{ ha} = 120\,000\text{ m}^2$.\\
\textbf{Đáp án của câu hỏi là A}
}
\end{ex}

\begin{ex} %[5D6H34]
Viết số đo diện tích $48\text{ km}^2\,9\text{ ha}$ dưới dạng số thập phân có đơn vị ki-lô-mét vuông là
\choice
{$48{,}9\text{ km}^2$}
{\True $48{,}09\text{ km}^2$}
{$48{,}009\text{ km}^2$}
{$489\text{ km}^2$}
\loigiai{
$48\text{ km}^2\,9\text{ ha} = 48\dfrac{9}{100}\text{ km}^2 = 48{,}09\text{ km}^2$.\\
\textbf{Đáp án của câu hỏi là B}
}
\end{ex}

\begin{ex} %[5D6H34]
Kết quả của phép tính $28{,}6\text{ ha} \times 5$ là
\choice
{$140{,}3\text{ ha}$}
{$14{,}3\text{ ha}$}
{$1\,430\text{ ha}$}
{\True $143\text{ ha}$}
\loigiai{
Ta có: $28{,}6 \times 5 = 143\text{ (ha)}$.\\
\textbf{Đáp án của câu hỏi là D}
}
\end{ex}

\begin{ex} %[5D6H34]
Đổi $\dfrac{5}{6}\text{ giờ}$ ra phút ta được kết quả là
\choice
{\True $50\text{ phút}$}
{$30\text{ phút}$}
{$45\text{ phút}$}
{$55\text{ phút}$}
\loigiai{
$\dfrac{5}{6}\text{ giờ} = \dfrac{5}{6} \times 60 = 50\text{ phút}$.\\
\textbf{Đáp án của câu hỏi là A}
}
\end{ex}

\begin{ex} %[5D6V34]
Một khu đất xây dựng nhà máy hình chữ nhật có chiều dài $300\text{ m}$, chiều rộng bằng $\dfrac{1}{3}$ chiều dài. Diện tích khu đất đó tính theo héc-ta là
\choice
{$30\text{ ha}$}
{\True $3\text{ ha}$}
{$0{,}3\text{ ha}$}
{$300\text{ ha}$}
\loigiai{
Chiều rộng khu đất là: $300 \times \dfrac{1}{3} = 100\text{ (m)}$.\\
Diện tích khu đất là: $300 \times 100 = 30\,000\text{ (m}^2\text{)} = 3\text{ ha}$.\\
\textbf{Đáp án của câu hỏi là B}
}
\end{ex}

\begin{ex} %[5D6V34]
Đường từ nhà Hùng đến trường gồm đoạn xuống dốc dài $650\text{ m}$ và đoạn lên dốc dài gấp đôi đoạn xuống dốc. Quãng đường từ nhà Hùng đến trường dài số ki-lô-mét là
\choice
{$1{,}3\text{ km}$}
{$19{,}5\text{ km}$}
{$1\,950\text{ km}$}
{\True $1{,}95\text{ km}$}
\loigiai{
Đoạn lên dốc dài là: $650 \times 2 = 1\,300\text{ (m)}$.\\
Cả quãng đường dài là: $650 + 1\,300 = 1\,950\text{ (m)} = 1{,}95\text{ km}$.\\
\textbf{Đáp án của câu hỏi là D}
}
\end{ex}

\begin{ex} %[5D6V34]
Một khu sinh thái có diện tích $128\text{ ha}$. Trong đó, $\dfrac{7}{8}$ diện tích là rừng đặc dụng, còn lại là khu vui chơi và dịch vụ. Diện tích khu vui chơi và dịch vụ tính theo ki-lô-mét vuông là
\choice
{$16\text{ km}^2$}
{$1{,}6\text{ km}^2$}
{\True $0{,}16\text{ km}^2$}
{$112\text{ ha}$}
\loigiai{
Diện tích khu vui chơi và dịch vụ chiếm: $1 - \dfrac{7}{8} = \dfrac{1}{8}$ (diện tích cả khu).\\
Diện tích khu vui chơi là: $128 \times \dfrac{1}{8} = 16\text{ (ha)} = 0{,}16\text{ km}^2$.\\
\textbf{Đáp án của câu hỏi là C}
}
\end{ex}

\begin{ex} %[5D6V34]
Đổi $\dfrac{4}{5}\text{ thế kỉ}$ ra năm ta được kết quả là
\choice
{$40\text{ năm}$}
{$20\text{ năm}$}
{$85\text{ năm}$}
{\True $80\text{ năm}$}
\loigiai{
Vì $1\text{ thế kỉ} = 100\text{ năm}$ nên $\dfrac{4}{5}\text{ thế kỉ} = \dfrac{4}{5} \times 100 = 80\text{ năm}$.\\
\textbf{Đáp án của câu hỏi là D}
}
\end{ex}

\begin{ex} %[5D6V34]
Một người đi xe máy khởi hành lúc $7\text{ giờ } 15\text{ phút}$ và đến nơi lúc $8\text{ giờ } 45\text{ phút}$. Vận tốc của người đó là $36\text{ km/h}$. Quãng đường người đó đã đi được là
\choice
{\True $54\text{ km}$}
{$36\text{ km}$}
{$45\text{ km}$}
{$72\text{ km}$}
\loigiai{
Thời gian đi là: $8\text{ giờ } 45\text{ phút} - 7\text{ giờ } 15\text{ phút} = 1\text{ giờ } 30\text{ phút} = 1{,}5\text{ giờ}$.\\
Quãng đường đi được là: $36 \times 1{,}5 = 54\text{ (km)}$.\\
\textbf{Đáp án của câu hỏi là A}
}
\end{ex}

\begin{ex} %[5D6V34]
Một xe tải chở $2$ tấn $400\text{ kg}$ gạo. Người ta dỡ bớt xuống $\dfrac{3}{8}$ số gạo đó. Số ki-lô-gam gạo còn lại trên xe là
\choice
{$900\text{ kg}$}
{$1\,800\text{ kg}$}
{\True $1\,500\text{ kg}$}
{$1\,200\text{ kg}$}
\loigiai{
Đổi $2\text{ tấn } 400\text{ kg} = 2\,400\text{ kg}$.\\
Số gạo dỡ xuống là: $2\,400 \times \dfrac{3}{8} = 900\text{ (kg)}$.\\
Số gạo còn lại là: $2\,400 - 900 = 1\,500\text{ (kg)}$.\\
\textbf{Đáp án của câu hỏi là C}
}
\end{ex}

\textbf{Phần 2. Trắc nghiệm đúng sai}
\setcounterpageref{ex}{0}

\begin{ex} %[5D6H34]
Xét tính đúng sai của các phép đổi đơn vị đo sau:
\choiceTF
{\True $5\text{ m } 360\text{ mm} = 5{,}36\text{ m}$.}
{$6\text{ tấn } 205\text{ kg} = 6{,}25\text{ tấn}$.}
{\True $6{,}5\text{ ha} = 65\,000\text{ m}^2$.}
{\True $\dfrac{7}{12}\text{ phút} = 35\text{ giây}$.}
\loigiai{
\begin{itemize}
    \item $5\text{ m } 360\text{ mm} = 5{,}36\text{ m}$.
    \item $6\text{ tấn } 205\text{ kg} = 6{,}205\text{ tấn} \neq 6{,}25\text{ tấn}$.
    \item $6{,}5\text{ ha} = 6{,}5 \times 10\,000 = 65\,000\text{ m}^2$.
    \item $\dfrac{7}{12}\text{ phút} = \dfrac{7}{12} \times 60 = 35\text{ giây}$.
\end{itemize}
\begin{itemize}
    \item a) \textbf{Đúng}.
    \item b) \textbf{Sai}: $6{,}205\text{ tấn}$.
    \item c) \textbf{Đúng}.
    \item d) \textbf{Đúng}.
\end{itemize}
}
\end{ex}

\begin{ex} %[5D6V34]
Khu sinh thái có diện tích $128\text{ ha}$. Trong đó $\dfrac{7}{8}$ diện tích là rừng đặc dụng, phần còn lại là khu vui chơi và dịch vụ. Xét tính đúng sai của các khẳng định sau:
\choiceTF
{\True Diện tích toàn bộ khu sinh thái bằng $1{,}28\text{ km}^2$.}
{\True Diện tích rừng đặc dụng là $112\text{ ha}$.}
{Diện tích khu vui chơi và dịch vụ là $16\text{ km}^2$.}
{\True Diện tích khu vui chơi và dịch vụ là $0{,}16\text{ km}^2$.}
\loigiai{
\begin{itemize}
    \item $128\text{ ha} = 1{,}28\text{ km}^2$.
    \item Rừng đặc dụng: $128 \times \dfrac{7}{8} = 112\text{ ha}$.
    \item Khu vui chơi: $128 - 112 = 16\text{ ha} = 0{,}16\text{ km}^2 \neq 16\text{ km}^2$.
\end{itemize}
\begin{itemize}
    \item a) \textbf{Đúng}.
    \item b) \textbf{Đúng}.
    \item c) \textbf{Sai}.
    \item d) \textbf{Đúng}.
\end{itemize}
}
\end{ex}

\textbf{Phần 3. Trả lời ngắn}
\setcounterpageref{ex}{0}

\begin{ex} %[5D6H34]
Đổi $93{,}17\text{ km}^2 : 7$ ra ki-lô-mét vuông dưới dạng số thập phân.
\shortans{13,31}
\loigiai{
\textbf{Đáp số:} $13{,}31$\\
$93{,}17 : 7 = 13{,}31\text{ (km}^2\text{)}$.
}
\end{ex}

\begin{ex} %[5D6H34]
Đổi $634\text{ cm}$ ra mét dưới dạng số thập phân.
\shortans{6,34}
\loigiai{
\textbf{Đáp số:} $6{,}34$\\
$634\text{ cm} = 6{,}34\text{ m}$.
}
\end{ex}

\begin{ex} %[5D6H34]
Khu đất hình chữ nhật $300\text{ m} \times 100\text{ m}$. Tính diện tích khu đất theo héc-ta.
\shortans{3}
\loigiai{
\textbf{Đáp số:} $3$\\
$300 \times 100 = 30\,000\text{ m}^2 = 3\text{ ha}$.
}
\end{ex}

\begin{ex} %[5D6V34]
Đường từ nhà Hùng đến trường gồm đoạn xuống dốc $650\text{ m}$ và lên dốc $1\,300\text{ m}$. Tính tổng chiều dài con đường theo ki-lô-mét. (Nhập kết quả dưới dạng số thập phân).
\shortans{1,95}
\loigiai{
\textbf{Đáp số:} $1{,}95$\\
$650 + 1\,300 = 1\,950\text{ m} = 1{,}95\text{ km}$.
}
\end{ex}

\begin{ex} %[5D6V34]
Khu sinh thái $128\text{ ha}$ có $\dfrac{7}{8}$ là rừng, còn lại là khu vui chơi. Tính diện tích khu vui chơi theo héc-ta.
\shortans{16}
\loigiai{
\textbf{Đáp số:} $16$\\
$128 \times \left(1 - \dfrac{7}{8}\right) = 16\text{ ha}$.
}
\end{ex}

\begin{ex} %[5D6V34]
Đổi $\dfrac{7}{12}\text{ phút}$ ra giây.
\shortans{35}
\loigiai{
\textbf{Đáp số:} $35$\\
$\dfrac{7}{12} \times 60 = 35\text{ (giây)}$.
}
\end{ex}
